% TOP_PROBLEM: {"id": 3106, "title": "Shannon capacity of the seven-cycle", "category_id": 3, "category": {"id": 3, "name": "graph_theory", "display_name": "Graph Theory", "description": "Problems involving graphs, networks, and their properties.", "slug": "graph-theory", "order_index": 3, "created_at": "2026-07-31T15:26:25.671Z"}, "status": "open", "problem_number": "OPG-36879", "set_id": 12, "statusQualification": "The conventional Shannon-capacity normalization is used; the source’s extra factor 1/k under the kth root does not change its limiting value."}
% FIRST_POSTED: 2026-10-01
% ENTRY_KIND: statement_only
% UnsolvedMath ID 3106; source catalogue code OPG-36879
% !TeX program = lualatex
% Definition and sources only; this file is not an LLM solution attempt.
\documentclass[11pt,a4paper]{article}
\usepackage[margin=25mm]{geometry}
\usepackage{fontspec}
\setmainfont{lmroman10-regular.otf}[BoldFont=lmroman10-bold.otf,
 ItalicFont=lmroman10-italic.otf,BoldItalicFont=lmroman10-bolditalic.otf]
\usepackage{amsmath,amssymb,mathtools}
\usepackage{unicode-math}
\setmathfont{latinmodern-math.otf}
\AtBeginDocument{\let\setminus\smallsetminus}
\usepackage{xurl,hyperref}
\hypersetup{colorlinks=true,urlcolor=blue}
\setcounter{secnumdepth}{0}
\setlength{\parindent}{0pt}
\setlength{\parskip}{5pt}
\setlength{\emergencystretch}{3em}
\begin{document}
\section*{Shannon capacity of the seven-cycle}\label{3106}
{\small UnsolvedMath ID 3106 · OPG-36879 · Graph Theory · OpenGarden}
\subsection{Definitions and mathematical statement}
Let $C_7$ be the simple cycle with vertices $\mathbb Z/7\mathbb Z$ and
edges between residues whose difference is $1$ or $-1$ modulo $7$.
For $k\ge1$, its $k$th strong power $C_7^{\boxtimes k}$ has vertex set
$(\mathbb Z/7\mathbb Z)^k$. Distinct tuples $x,y$ are adjacent if,
in every coordinate, they are equal or adjacent in $C_7$.

For a finite graph $H$, the independence number $\alpha(H)$ is the
largest cardinality of a vertex set containing no adjacent pair.
The Shannon capacity is
\[
 \Theta(C_7)=\lim_{k\to\infty}
       \alpha(C_7^{\boxtimes k})^{1/k}
       =\sup_{k\ge1}\alpha(C_7^{\boxtimes k})^{1/k}.
\]
The limit exists because independent sets in two powers can be multiplied
to give an independent set in their combined power, making their
independence numbers supermultiplicative. All quantities are finite and
bounded after taking $k$th roots by seven.

Determine this real number exactly. In coding language, a set of tuples
is independent when any two distinct words have some coordinate whose
symbols are neither equal nor consecutive on the seven-cycle.
The capacity is the asymptotically achievable number of distinguishable
symbols per coordinate, not its logarithm.

The source discussion includes a factor $1/k$ inside its root; this
factor has limiting root one and does not alter the limit. The conventional
normalization above makes the supremum formula transparent. Computing one
finite strong power supplies a lower bound but does not by itself
determine the limit. The notation $\Theta$ distinguishes capacity from
other graph parameters sometimes denoted by a lowercase theta.
\subsection{Short English statement}
What is the exact exponential growth rate of the largest independent sets in repeated strong products of the seven-cycle?
\subsection{Sources}
\begin{itemize}
\item[S1] ULAM.AI and the UnsolvedMath Contributors, supplied problems.json, record 3106 (OPG-36879), OpenGarden collection. Catalogue identity and source transcription; CC BY 4.0.
\url{https://huggingface.co/datasets/ulamai/UnsolvedMath}.
\item[S2] Open Problem Garden, Shannon capacity of the seven-cycle. Original problem and discussion; the mathematical formulation below is expanded from the supplied source record.
\url{http://www.openproblemgarden.org/op/shannon_capacity_of_the_seven_cycle}.
\end{itemize}
\end{document}
